\documentclass[a4paper, 12pt]{article}
\usepackage{a4wide}
\usepackage {amsmath}
\usepackage{amssymb}
\usepackage {graphicx}
\usepackage[utf8]{inputenc} 
\usepackage[francais]{babel}
\usepackage{fancyhdr}
\usepackage{setspace}
\usepackage{fancyhdr}
\usepackage{lastpage}
\usepackage{extramarks}
\usepackage{chngpage}
\usepackage{soul}
\usepackage{algorithmicx} 
\usepackage{algpseudocode} 
\usepackage{multicol}
\usepackage[usenames,dvipsnames]{color}
\usepackage{graphicx,float,wrapfig}
\usepackage{ifthen}
\usepackage{listings}
\usepackage{courier}
\usepackage{esint}
\usepackage{bbm}
\usepackage{graphics}
\usepackage{graphicx}
\usepackage{subfig}
\usepackage{epsfig}
\usepackage{pgf,tikz}
\usetikzlibrary{arrows}
\usepackage{braket}
\usepackage{MnSymbol,wasysym}
\usepackage{marvosym}
\usepackage{dsfont}


\lhead{} 
\chead{} 
\rhead{\bfseries Matlab} 
\lfoot{Q \& JC}
%\cfoot{} 
%\rfoot{\thepage}

% This is the color used for MATLAB comments below
\definecolor{MyDarkGreen}{rgb}{0.0,0.4,0.0}

% For faster processing, load Matlab syntax for listings
\lstloadlanguages{Matlab}%
\lstset{language=Matlab,                        % Use MATLAB
        frame=single,                           % Single frame around code
        basicstyle=\small\ttfamily,             % Use small true type font
        keywordstyle=[1]\color{Blue}\bf,        % MATLAB functions bold and blue
        keywordstyle=[2]\color{Purple},         % MATLAB function arguments purple
        keywordstyle=[3]\color{Blue}\underbar,  % User functions underlined and blue
        identifierstyle=,                       % Nothing special about identifiers
                                                % Comments small dark green courier
        commentstyle=\usefont{T1}{pcr}{m}{sl}\color{MyDarkGreen}\small,
        stringstyle=\color{Purple},             % Strings are purple
        showstringspaces=false,                 % Don't put marks in string spaces
        tabsize=5,                              % 5 spaces per tab
        %
        %%% Put standard MATLAB functions not included in the default
        %%% language here
        morekeywords={xlim,ylim,var,alpha,factorial,poissrnd,normpdf,normcdf},
        %
        %%% Put MATLAB function parameters here
        morekeywords=[2]{on, off, interp},
        %
        %%% Put user defined functions here
        morekeywords=[3]{FindESS, homework_example},
        %
        morecomment=[l][\color{Blue}]{...},     % Line continuation (...) like blue comment
        numbers=left,                           % Line numbers on left
        firstnumber=1,                          % Line numbers start with line 1
        numberstyle=\tiny\color{Blue},          % Line numbers are blue
        stepnumber=5                            % Line numbers go in steps of 5
        }

% Includes a MATLAB script.
% The first parameter is the label, which also is the name of the script
%   without the .m.
% The second parameter is the optional caption.
\newcommand{\matlabscript}[2]
  {\begin{itemize}\item[]\lstinputlisting[caption=#2,label=#1]{#1.m}\end{itemize}}

\pagestyle{fancy}

\begin{document}

\bibliographystyle{alpha}

\title{Bases en programmation Matlab}

\author{Quentin Rafhay - Jean-Christophe Toussaint\\
%  Phelma\\
%\texttt{jean-christophe.toussaint@phelma.grenoble-inp.fr}
}
\date{\today}
 
\maketitle


L'ensemble des exercices, fourni ci-après, est à réaliser en totale autonomie et constitue le socle
de base en Matlab que vous devez assimiler.

\section{Exercice : multiplication matricielle}

Matlab est un langage adapté à l'algèbre linéaire et aux calculs matriciels.

\begin{enumerate} 
 \item En utilisant toutes les fonctionnalités de Matlab, développer une fonction mult(A,B)
retournant le produit matricielle $R=A\;B$. 
\item Pour aider l'utilisateur, compléter la fonction précédente avec un commentaire
en début du fichier mult.m, donnant un mode d'emploi sommaire.
\item Pour professionnaliser votre code, insérer un bloc permettant de vérifier la compatibilité
des tailles des matrices. Utiliser les fonctions {\tt size} et {\tt error}. Cette dernière permet stopper 
le déroulement de votre programme.
\item Développer une nouvelle fonction  mult2(A, B) retournant le produit matriciel $R=A\;B$, dans laquelle
l'opération est faite de manière explicite sous la forme de boucles sur les éléments des matrices.
\item En plaçant les fonctions {\tt tic} et {\tt toc} respectivement, avant et après l'appel d'une fonction, Matlab
retourne le temps de calcul CPU. Comparer l'éfficacité des différentes implémentations proposées
en vous limitant aux matrices carrées.
Utiliser la fonction {\tt rand} pour générer des matrices de grande taille. 
\end{enumerate} 

\section{Exercice : manipulation de tableaux}

On considère un tableau $A$ unidimensionnel de taille $N=1000$ voire plus. 
\begin{enumerate} 
\item générer le tableau (ligne) $A$ avec la fonction Matlab {\tt rand}.
\item Développer une fonction min\_max(A) retournant les valeurs minimale et maximale.

\item Dévélopper une fonction tri\_bulle(A) permettant de trier le tableau de manière croissante.
La méthode la plus simple dite de "tri par bulle" consiste à permuter si nécessaire, deux éléments successifs
du tableau et à recommencer dès le début du tableau tant qu'il n'est pas trié.

\item  Hoare a inventé en 1962, une méthode de type "diviser pour régner"  permettant de trier un
tableau en $N \ln N$ opérations. Son algorithme s'écrit :

\begin{algorithmic}[1]
   \begin{multicols}{2}
\Function{Tri\_rapide}{$T,g,d$}
    \If {$g < d$}
        \State $ip, T \gets \Call{Partition}{$T, g, d$} $
         \State $T \gets \Call{Tri\_rapide}{$T, g, ip$} $
         \State $T \gets \Call{Tri\_rapide}{$T, ip+1, d$} $
    \EndIf
    \State \Return $T$
\EndFunction
\\
\Function{Partition}{$T,g,d$}
    \State $pivot \gets T[g]$
    \State $i \gets g-1$
    \State $j \gets d+1$
        \columnbreak
    \While{$true$}
    \Repeat
    \State $j \gets j-1$
    \Until{$T[j] <= pivot$}

    \Repeat
    \State $i \gets i+1$
    \Until{$T[i] >= pivot$}
    \If {$i < j$}
    \State $T \gets \Call{Echanger}{$T, i, j$} $
    \Else
    \State \Return $j, T$
    \EndIf
        \EndWhile
\EndFunction
\end{multicols}
\end{algorithmic}


\item La boucle conditionnelle {\tt repeat ... until (condition)} n'existe pas en Matlab, proposez
une solution de remplacement. N'oubliez pas de la tester à part.

\item Implémenter l'algorithme suivant sous la forme d'une fonction Matlab. Son appel 
A=tri\_rapide(A, 1, length(A)) retourne le tableau trié.


\item Faites un programme appelant la fonction tri\_rapide permettant de trier des tableaux
de taille $2^n$ où $n \in [1, 20]$.  En utilisant les fonctions {\tt tic} et {\tt toc} , stocker le temps CPU
dans un tableau et finalement tracer l'évolution du temps de calcul en fonction de la taille.

\item Vérifier qu'il varie en $N \ln N$ dans la limite des grandes valeurs de $N$.

%\begin{figure}[h]
%   \caption{\label{Algo} Algorithme Tri Rapide}
%   \includegraphics[width=8cm]{Algo_Tri_Rapide}
%\end{figure}

\end{enumerate} 

\section{Exercice : nombres premiers}

On désire connaitre la liste $Lp$ des nombres premiers inférieurs à $N$ entier donné.
\begin{enumerate} 
\item Développer une fonction Matlab retournant cette liste.\\

{\it Indication } : on peut construire la liste de manière dynamique sans allocation préalable.
L'initialisation de la liste vide se fait par l'instruction {\tt Lp=[ ]} et l'insertion
à droite d'un élément $e$ dans $Lp$ par {\tt Lp = [Lp, e]}.
\end{enumerate} 
 
 \section{Méthode de la puissance itérée}

\subsection{Notations}

On considère une matrice réelle symétrique ou hermitique $A$ de dimensions $n \times n$.
Elle vérifie par conséquent la propriété $A=(A^t)^*$ où $*$ est l'opérateur de conjugaison
complexe que l'on applique à chaque élément de $A^t$. On sait aussi que dans ce cas, $A$ est toujours diagonalisable dans $\mathbbm C$.

On définit le produit scalaire de deux vecteurs  $| u\rangle=\left(x_1, x_2, \cdots, x_n \right)^t$ et $| v\rangle=\left(y_1, y_2, \cdots, y_n \right)^t$ par:
\begin{equation}
\langle v | u \rangle= \sum_{i=1}^{n} x_i \; y_i^*
\end{equation}

Supposons pour simplifier que le spectre de ses valeurs propres est non dégénéré.
Autrement dit, à chaque valeur propre $\lambda_i$ correspond un s.e.v. de dimension 1,
décrit par un vecteur propre $| u_i\rangle$.  La base formée par l'ensemble des vecteurs propres
est orthogonale et peut être normée.


On peut décomposer $A$ comme une somme de projecteurs :

\begin{equation}
A= \sum_{i=1}^{n} \lambda_i | u_i\rangle \langle u_i |
\end{equation}

\subsection{Recherche de la plus grande valeur propre en module}

On désire dans un premier temps extraire la plus grande valeur propre en module que nous nommons 
$\lambda_1$ et définir le
vecteur propre associé $| u_1\rangle$ en utilisant une méthode itérative.

Décrivons cet algorithme.
On choisit d'abord un vecteur aléatoire $|x_0 \rangle$ de dimension $n$ sur lequel on applique $p$
fois la matrice $A$. 

A la première itération, on a :
\begin{equation}
|x_1 \rangle = A |x_0 \rangle = \sum_{i=1}^{n} \lambda_i | u_i\rangle \langle u_i |x_0 \rangle
\end{equation}

A l'itération $p$, on obtient :
\begin{equation}
|x_p \rangle = A^p |x_0 \rangle = \sum_{i=1}^{n} \lambda_i^p | u_i\rangle \langle u_i |x_0 \rangle
\end{equation}

que l'on peut développer suivant la forme suivante:
\begin{equation}
|x_p \rangle = A^p |x_0 \rangle =\lambda_1^p \left( \langle u_1 |x_0 \rangle \, | u_1\rangle 
+\left(\frac{\lambda_2}{\lambda_1}\right)^p \langle u_2 |x_0 \rangle \, | u_2\rangle +
\cdots
+\left(\frac{\lambda_n}{\lambda_1}\right)^p \langle u_n |x_0 \rangle \, | u_n\rangle 
\right) \end{equation}

En remarquant que $|\frac{\lambda_i}{\lambda_1}| \leq 1$ et dans la limite des grandes valeurs
de $p$, on observe que:
\begin{equation}
\lim_{p \rightarrow \infty}
|x_p \rangle \propto \langle u_1 |x_0 \rangle \, | u_1\rangle  
\end{equation}

Le vecteur $|x_p \rangle$ tend à s'aligner vers le vecteur propre $|u_1 \rangle$.
\subsection{Implémentation de la méthode en Matlab/Octave}

\begin{enumerate}
\item On donne ci-dessous un algorithme simplifié mettant en oeuvre la méthode de la puissance itérée.
%\insertlisting{puissance.m}{puissance}


\begin{algorithmic}[1]
\begin{multicols}{2}
       % \columnbreak
\Function{puissance}{$A, x$}
     \State $iter \gets 0$   
\State $\epsilon \gets 10^{-6} $
        \Repeat
    \State $x0 \gets x$
    \State $v \gets A \; x $
    \State $x \gets v /  \Call{norm}{v} $  

     \State $iter \gets iter+1$   
         \Until{$\Call{norm}{x-x0} < \epsilon$}
 \\           
\State $\lambda \gets x^t\; A\;  x$ 
    \State \Return $\lambda, x$
\EndFunction
\end{multicols}
\end{algorithmic}


\item Implémenter l'algorithme précédent sous la forme d'une fonction Matlab {\tt [lb, u] =puissance(A, x)} où A est une matrice de dimensions $n \times n$ et un vecteur colonne $x$ 
non nul, de dimension $n$.
Elle retourne deux variables 
{\tt lb} et {\tt u} qui correspondent à la valeur propre $\lambda_1$ et au vecteur propre $|u_1 \rangle$.

\item  Compléter le fichier $\tt{puissance.m} \rm$ avec un en-tête de commentaires résumant 
l'usage de la fonction. 

\item Tester son comportement depuis la fenêtre de commande, en prenant 
$A =
\left (
\begin{array}{c c c} 1 & 2 & 3 \\ 2 & 1 & 2 \\ 3 & 2 & 1 \end{array}
\right ) 
$.

\item Vérifier que les résultats sont cohérents avec ceux donnés par la fonction interne à
Matlab $\tt{eig(A)} \rm$ et $\tt{[V, D] = eig(A)} \rm$. Commenter.
\item Modifier la fonction pour recupérer le nombre d'itérations. Quel est l'intérêt de le faire?
\item Exécuter de nouveau cette fonction en prenant comme matrice $-A$.  Commenter.

\item Que se passe-t-il si la plus grande valeur propre en module est dégénérée? 
La méthode fonctionne-t-elle encore dans ce cas?

\item On peut générer une matrice aléatoire de grande taille avec la fonction {\tt rand(N, N)}
dont la valeur de chaque élément est dans $[0, 1[$. Comment la rendre symétrique?

\item Est-il raisonnable d'utiliser cet algorithme dans le cas d'une matrice non symétrique?

\item Est-il raisonnable d'utiliser cet algorithme dans le cas d'une matrice hermitique?

\item Modifier la fonction  {\tt [l, u] =puissance(A, x)} en imposant
un nombre maximal d'itérations.  Quelles en sont les conséquences immédiates?

\item Proposez une approche, même si elle n'est pas efficace, pour extraire la plus petite valeur propre en module d'une matrice.

\item Montrer qu'en modifiant la ligne 6 de l'algorithme précédent, la plus petite valeur propre $\lambda_{min}$
et son vecteur propre associé peuvent être calculés à condition touetfois que $\lambda_{min}>0$.

\end{enumerate} 
\end{document}

